\section{Einleitung} \label{einleitung}

Dieses Kapitel führt in die Motivation, den Lösungsansatz, die Zielsetzung und die Abgrenzung der Projektarbeit ein.

\subsection{Herausforderungen von Reverse Engineering und Re-Implementierung}

Das Reverse Engineering und die Re-Implementierung bestehender Codebasen zählen zu den zentralen Herausforderungen der zeitgenössischen Softwareentwicklung. Im Lebenszyklus gewachsener Systeme entstehen technologisch bedingte Grenzen hinsichtlich Performance, Wartbarkeit und Skalierbarkeit, sodass ein Wechsel auf moderne Sprachen und Architekturen erhebliche Vorteile verspricht. Dazu zählen insbesondere verbesserte Typsicherheit, höhere Ressourceneffizienz und nativer Concurrency-Support. Sowohl klassische modellgetriebene Arbeiten als auch aktuelle Übersichten zu LLM-gestütztem Reverse Engineering ordnen diese Transformationsaufgabe als besonders anspruchsvoll ein \parencite{raibulet2017mdre,hu2025sokreverseengineering}.

In der Praxis scheitert eine solche Migration jedoch häufig an der hohen kognitiven Last für Entwickler, einer lückenhaften Dokumentation des Bestands sowie an grundlegenden Paradigmenwechseln zwischen Quell- und Zieltechnologie. Gerade die Überführung bestehender Systeme in explizite, rekonstruierbare Anforderungen bildet daher eine zentrale Brücke zwischen Reverse Engineering und späterer Neuimplementierung \parencite{fahmi2007reverse}.

\subsection{Spec-Driven Development durch Agentenorchestrierung}

Als Lösungsansatz wird eine Migrationsstrategie verfolgt, die nicht auf direkter Transpilierung basiert, sondern auf einer vorgelagerten Spezifikationsebene. Im Sinne von Spec-Driven Development (SDD) wird der bestehende Quellcode zunächst in technische Anforderungen und Verhaltensbeschreibungen überführt, bevor die Implementierung in der Zielsprache erfolgt. Eine solche explizite Spezifikation ist auch deshalb vielversprechend, weil präzise Angaben zu I/O-Formaten, Bedingungen und Beispielen die Zuverlässigkeit nachgelagerter Codegenerierung erhöhen \parencite{midolo2026promptguidelines}.

Die Erstellung, Prüfung und Nutzung dieser Spezifikationen soll durch einen orchestrierten Verbund spezialisierter Agenten erfolgen. Dadurch werden Analyse, Spezifikation, Verifikation und Generierung in klar abgegrenzte Rollen zerlegt. Vergleichbare arbeitsteilige Muster wurden bereits in agentischen Software-Engineering-Ansätzen untersucht \parencite{qian2024chatdev,yang2024sweagent}.

\subsection{Zielsetzung der Migrations-Pipeline}

Ziel der Projektarbeit ist die Konzeption und prototypische Implementierung einer automatisierten Migrations-Pipeline auf Basis eines Multi-Agenten-Systems (MAS). Der Prototyp soll zeigen, dass sich aus bestehendem Quellcode präzise technische Spezifikationen in Markdown gewinnen lassen, die anschließend als belastbare Grundlage für die Code-Generierung dienen. Die technische Kopplung zwischen Agenten und Analysewerkzeugen wird dabei bewusst über standardisierte Schnittstellen gedacht \parencite{mcpdocs}.

Der angestrebte Proof of Concept konzentriert sich auf ein repräsentatives Quell-Modul. Damit soll ein vollständiger Durchstich durch Analyse, Spezifikation, Verifikation und Generierung demonstriert werden, ohne den Rahmen der Projektarbeit zu überschreiten.

\subsection{Abgrenzung auf textbasierte Spezifikationen}

Die Projektarbeit fokussiert auf textbasierte technische Spezifikationen als zentrales Zwischenartefakt. Visuelle Modelle stehen noch nicht im Zentrum des prototypischen Umsetzungsumfangs. Diese Abgrenzung erlaubt es, zunächst die Qualität textueller Anforderungen, ihrer Verifikation und ihrer Wirkung auf die nachgelagerte Codeerzeugung in den Vordergrund zu stellen \parencite{midolo2026promptguidelines,korn2026reportingprompting}.

Nicht Ziel ist die vollständige Automatisierung beliebiger Migrationsszenarien, sondern die nachvollziehbare Demonstration eines Kernprozesses innerhalb eines überschaubaren Anwendungsfalls.
